\chapter{Modelos de redes}

\section{Estructura comun}
Un modelo de redes representa entidades como nodos y relaciones como arcos. Esta
abstraccion cubre rutas mas cortas, arboles de expansion minima, flujo maximo,
transporte, asignacion y vendedor viajero \cite{ahuja1993network}.

\section{Familias contempladas en QSB}
El formato normalizado de \qsb{} usa un sobre JSON con un campo
\texttt{kind} y un objeto \texttt{model}. Las familias documentadas actualmente
incluyen:
\begin{longtable}{ll}
\toprule
Codigo & Familia \\
\midrule
SPP & Ruta mas corta \\
MST & Arbol de expansion minima \\
MFP & Flujo maximo \\
TSP & Vendedor viajero \\
AP & Asignacion \\
TP & Transporte \\
\bottomrule
\end{longtable}

\section{Criterio editorial para nuevos casos}
Cada caso de red debera incluir diagrama, matriz o lista de arcos, formulacion
matematica, archivo JSON de entrada, comando de solucion y verificacion manual
del costo, flujo o asignacion resultante.
